\documentclass[12pt]{article}
\usepackage[utf8]{inputenc}
\usepackage[T1]{fontenc}
\usepackage[margin=1in]{geometry}
\usepackage{setspace}
\usepackage{xcolor}
\usepackage[colorlinks=true, urlcolor=blue!60!black, linkcolor=blue!60!black]{hyperref}
\setlength{\parindent}{0em}
\setlength{\parskip}{0.45em}
\renewcommand{\baselinestretch}{1.0}
\begin{document}
\begin{center}
{\Large \textit{Wickard v. Filburn}, 317 U.S. 111 (1942)\par}
\vspace{0.3em}
{\normalsize Case Brief --- One-Page Edition\par}
\end{center}

\textbf{Memory Jogger.} The home-grown wheat case: substantial effects and aggregation (\href{https://supreme.justia.com/cases/federal/us/317/111/}{317 U.S. 111}).

\textbf{Facts.} Under the Agricultural Adjustment Act's wheat quotas, Ohio farmer Roscoe Filburn was allotted 11.1 acres but sowed 23, using much of the crop on his farm. He refused to pay a \$117.11 penalty on his excess wheat and sued Secretary Wickard.

\textbf{Procedural History.} A three-judge district court enjoined the penalty; direct appeal.

\textbf{Votes.} Reversed, 9--0. Jackson wrote for a unanimous Court.

\textbf{Holding.} Local, non-commercial activity may be regulated if it substantially affects interstate commerce, and a trivial individual contribution is aggregated with that of others similarly situated. The penalty did not violate due process.

\textbf{Reasoning.} Formal labels like production and indirect effects give way to economic reality. Home-grown wheat overhangs the market and displaces purchases, so it competes with wheat in commerce; summed across many farmers, its effect is substantial. Regulating subsidized farmers is not arbitrary.

\textbf{Examples.} \emph{Same-Side:} home-grown marijuana (\textit{Raich}). \emph{Opposite-Side:} guns near schools (\textit{Lopez}). \emph{Fence-Sitter:} home distilling.

\textbf{Critique.} Originalists say \textit{Wickard} erased enumerated limits; pragmatists call it essential to a national economy.

\textbf{Key Quotation.} ``But even if appellee's activity be local, and though it may not be regarded as commerce, it may still, whatever its nature, be reached by Congress if it exerts a substantial economic effect on interstate commerce, and this irrespective of whether such effect is what might at some earlier time have been defined as `direct' or `indirect.' '' --- Justice Jackson, Opinion of the Court, 317 U.S. at 125.

\end{document}
